% Source: https://artofproblemsolving.com/community/c3h1569742_interesting_coincidence
% Stable identity: aops-topic-1569742-q01
Let $n=10a+b$ with digits $a,b\in\{1,\dots,9\}$, and let $r=10b+a$ be its reversal.

\textbf{Why some numbers work.}
\[
n^2=100a^2+10\cdot2ab+b^2,\qquad r^2=100b^2+10\cdot2ab+a^2 .
\]
If $a^2$, $2ab$ and $b^2$ are all at most $9$, there are no carries: the digits of $n^2$ are $a^2,\,2ab,\,b^2$ and those of $r^2$ are $b^2,\,2ab,\,a^2$, the same digits reversed. This happens for $(a,b)=(1,1),(1,2),(2,1),(1,3),(3,1),(2,2)$, i.e.\ $n=11,12,21,13,31,22$.

\textbf{No others.} Since $b\ne0$, $n^2$ does not end in $0$, so its reversal has as many digits as $n^2$. Hence $r^2$ and $n^2$ must have the same number of digits.
\begin{itemize}
\item \emph{Both have three digits.} Then $n\le31$ and $r\le31$, so $a,b\le3$, and $n,r\notin\{23,32,33\}$. This leaves exactly $11,12,13,21,22,31$, all of which work.
\item \emph{Both have four digits.} Then $n,r\ge32$. The last digit of $r^2$ is the last digit of $a^2$, and it must equal the first digit of $n^2$. But for $32\le n\le99$ the first digit of $n^2$ lies in a short range determined by $a$:
\[
\begin{array}{c|ccccccc}
a & 3 & 4 & 5 & 6 & 7 & 8 & 9\\ \hline
\text{first digit of }n^2 & 1 & 1,2 & 2,3 & 3,4 & 4,5,6 & 6,7 & 8,9\\
\text{last digit of }a^2 & 9 & 6 & 5 & 6 & 9 & 4 & 1
\end{array}
\]
(for example $a=7$: $70^2=4900\le n^2\le79^2=6241$). They never agree, so there is no four-digit case.
\end{itemize}
Therefore the numbers are exactly
\[
\boxed{11,\ 12,\ 13,\ 21,\ 22,\ 31},
\]
precisely those whose squaring involves no carries. (With a units digit of $0$, the numbers $10$, $20$, $30$ also ``work'' if leading zeros are dropped; the condition on the units digit excludes them.)
